\section{Alternative Data and Natural Language Processing}\label{sec:alternative_data}

Alternative data---non-traditional data sources including news, social media, satellite imagery, web traffic, and corporate filings---represents the fastest-growing application area for ML in finance. NLP methods are the primary tool for extracting signals from textual alternative data.

\subsection{Financial Sentiment Analysis}

\subsubsection{Dictionary-Based Methods}

The foundation of financial text analysis is the \citet{Loughran2011} dictionary, which demonstrated that general-purpose sentiment dictionaries (e.g., Harvard General Inquirer) perform poorly on financial text because words like ``liability,'' ``tax,'' and ``cost'' have neutral meanings in financial context. The Loughran-McDonald dictionary remains a benchmark against which ML approaches are evaluated.

\subsubsection{Pre-trained Language Models}

The transformer revolution has reached financial NLP:

\begin{itemize}[leftmargin=*]
    \item \textbf{FinBERT} \citep{Araci2019}: BERT fine-tuned on financial news. Achieves substantially higher accuracy than dictionary methods and general-purpose BERT on financial sentiment classification.

    \item \textbf{BloombergGPT} \citep{Wu2023}: A 50-billion parameter LLM trained on Bloomberg's proprietary corpus of financial documents. Outperforms general-purpose LLMs on financial NLP benchmarks while maintaining competitive general capabilities.

    \item \textbf{General LLMs (GPT-4, Claude):} \citet{LopezLira2023} show that GPT-3.5/4 can predict stock returns from news headlines, with performance comparable to purpose-trained models. The zero-shot capability eliminates the need for labeled training data.

    \item \textbf{Open-source finance LLMs:} FinGPT \citep{Yang2023} and InvestLM provide open-source alternatives for financial NLP, enabling reproducible research.
\end{itemize}

\subsection{News and Event Analysis}

Beyond sentiment, ML methods extract structured information from news:

\begin{itemize}[leftmargin=*]
    \item \textbf{Event detection:} Identifying market-moving events (M\&A announcements, earnings surprises, regulatory changes) from news feeds in real-time \citep{EinDor2020}.

    \item \textbf{Causal extraction:} Extracting causal relationships (``tariffs caused a decline in imports'') from financial news using transformer-based information extraction \citep{Mariko2022}.

    \item \textbf{Earnings call analysis:} ML models analyze the text and audio of earnings calls, extracting CEO tone, uncertainty language, and deceptive patterns that predict post-call stock movements \citep{Li2021, Qin2019}.
\end{itemize}

\subsection{Social Media and Crowd Signals}

\begin{itemize}[leftmargin=*]
    \item \textbf{Twitter/X:} \citet{Bollen2011} pioneered Twitter mood prediction of DJIA movements. More recent work uses transformer models on high-frequency tweet streams, with mixed evidence of genuine predictive content versus noise \citep{Renault2020}.

    \item \textbf{Reddit (WallStreetBets):} The GameStop episode (January 2021) catalyzed academic interest in retail investor coordination on social media. \citet{Bradley2021} analyze the information content of Reddit posts for stock prediction.

    \item \textbf{StockTwits:} A dedicated financial social platform that provides labeled bullish/bearish sentiment, making it a convenient training dataset for sentiment models \citep{Oliveira2017}.
\end{itemize}

\subsection{Non-Textual Alternative Data}

\begin{itemize}[leftmargin=*]
    \item \textbf{Satellite imagery:} CNNs applied to satellite images of parking lots (retail foot traffic), oil storage tanks (inventory levels), and agricultural fields (crop yields). \citet{Katona2021} demonstrate that satellite-derived parking lot occupancy predicts retailer revenue ahead of official reports.

    \item \textbf{Web data:} Job postings, product reviews, app download statistics, and web traffic as leading indicators of corporate performance \citep{Green2019}.

    \item \textbf{ESG and climate:} NLP extraction of ESG metrics from corporate reports, news, and regulatory filings. This subfield is growing rapidly due to regulatory pressure for ESG disclosure \citep{Kalesnik2022}.
\end{itemize}

\subsection{Challenges and Limitations}

\begin{enumerate}[leftmargin=*]
    \item \textbf{Signal decay:} Alternative data signals are rapidly incorporated into prices once widely adopted, creating an arms race for novel data sources.
    \item \textbf{Backtest inflation:} Alternative data is often available only for recent periods, limiting out-of-sample evaluation.
    \item \textbf{Labeling ambiguity:} Financial sentiment is context-dependent (``interest rates rose'' is positive for banks, negative for borrowers), making supervised training challenging.
    \item \textbf{Cost:} Commercial alternative data is expensive, creating barriers to academic research and replication.
\end{enumerate}
